\section{Functionals and Euler--Lagrange Equations}

\begin{definition}[Differentiability Classes]
    For \(\Omega \subseteq \RR^n\), we define
    \begin{align*}
        \contClass{k} \pqty{\Omega} &= \set{\functiondomain{y}{\Omega}{\RR}}{y \text{ is }k\text{ times continuously differentiable}}, \\
        \smoothClass \pqty{\Omega} &= \set{\functiondomain{y}{\Omega}{\RR}}{y \text{ is smooth (infinitely differentiable)}}
    \end{align*}
\end{definition}

For our discussion, usually \(\contClass{2}\) is sufficient. Denote \(\contClassSymb\) be some set of functions, for example \(\contClass{n} \pqty{\Omega}\).

\begin{definition}[Functional]
    A \emph{functional} is a map
    \[
        \function{F}{\contClassSymb}{\RR}{y}{F\bqty{y},}
    \]
    i.e. \(F\) maps a function to a number. More generally, a functional can take inputs of the form \(\pqty{f_1, \ldots, f_n}\) and output a number.
\end{definition}

An important case is when \(\Omega = \ccintv{\alpha}{\beta} \subseteq \RR\), and \(\contClassSymb = \smoothClass \pqty{\Omega}\), and
\[
    F\bqty{y} = \int_{\alpha}^{\beta} f\pqty{y \pqty{x}, y' \pqty{x}, x} \dd{x}
\]
for some smooth function \(f\) (with three arguments). We could generalise \(f\) to depend on other derivatives, but this is the most important example.

Our question is, which \(y \in \contClassSymb\) optimises \(F\)?

Consider varying \(y\) to \(y + \var{y} \in \contClassSymb\), and we have
\begin{align*}
    \var{F\bqty{y}} &= F\bqty{y + \var{y}} - F\bqty{y} \\
    &= \int_{\alpha}^{\beta} \bqty{f \pqty{y + \var{y}, y' + \var{y}, x} - f \pqty{y, y', x}} \dd{x} \\
    &= \int_{\alpha}^{\beta} \pqty{\var{y} \pdv{f}{y} + \var{y'} \pdv{f}{y'}} + \cdots\\
    &= \int_{\alpha}^{\beta} \dd{x} \var{y\pqty{x}} \pqty{\pdv{f}{y} - \dv{x} \pdv{f}{y'}} + \bqty{\var{y} \pdv{f}{y'}}_{\alpha}^{\beta} + \cdots
\end{align*}
using integration by parts.

Now assume that \(\contClassSymb\) is such that the boundary conditions makes the final term vanish, for example,
\begin{itemize}
    \item \(y\pqty{\alpha} = a, y\pqty{\beta} = b\) for some \(a, b\); or
    \item \(\pdv*{f}{y'} = 0\) at \(x = \alpha, \beta\), for example \(f = \sqrt{1 + y'^2}\) equipped with \(y'\pqty{\alpha} = y'\pqty{\beta} = 0\).
\end{itemize}

\begin{definition}[Functional Derivative]
    \label{def:functional-derivative}%
    We define
    \[
        \fdv{F\bqty{y}}{y\pqty{x}} = \pdv{f}{y} - \dv{x} \pdv{f}{y'}
    \]
    to be the \emph{functional derivative} of \(F\bqty{y}\) w.r.t \(y\pqty{x}\). This is a function in \(x\).
\end{definition}

So \(\var{F\bqty{y}}\) vanishes to first order for \emph{arbitrary} \(\var{y\pqty{x}}\), if and only if \(y\pqty{x}\) satisfies
\begin{equation}
    \fdv{F\bqty{y}}{y\pqty{x}} = \pdv{f}{y} - \dv{x} \pdv{f}{y'} = 0,
    \label{eq:euler-lagrange}
\end{equation}
the \emph{Euler--Lagrange Equation}. This is an second order ODE for \(y\pqty{x}\).

\begin{theorem}[Euler--Lagrange Equation]
    Let \(\functiondomain{f}{\RR^3}{\RR}\) be a function, and \(\contClassSymb\) be smooth on \(\ccintv{\alpha}{\beta}\) such that the boundary equation
    \[
        \bqty{\var{y} \pdv{f}{y'}}_{\alpha}^{\beta} = 0
    \]
    holds for all \(\var{y}, y \in \contClassSymb\). Then a function \(y \in \contClassSymb\) minimises the functional
    \[
        F\bqty{y} = \int_{\alpha}^{\beta} f\pqty{y, y', x} \dd{x}
    \]
    if and only if \autoref{eq:euler-lagrange} (the Euler--Lagrange Equation) holds.
\end{theorem}

For a generalisation, for a function
\[
    \vb{y} \pqty{x} = \pqty{y_1 \pqty{x}, y_2 \pqty{x}, \ldots, y_n \pqty{x}}
\]
for \(x \in \ccintv{\alpha}{\beta}\), i.e., \(\functiondomain{\vb{y}}{\ccintv{\alpha}{\beta}}{\RR^n}\), we define
\[
    F\bqty{\vb{y}} = \int_{\alpha}^{\beta} f\pqty{\vb{y}, \vb{y}', x} \dd{x}
\]
and we have
\[
    \var{F\bqty{y}} = \int_{\alpha}^{\beta} \sum_{i = 1}^{n} \var{y_i \pqty{x}} \pqty{\pdv{f}{y_i} - \dv{x} \pdv{f}{y'_i}} \dd{x} + \bqty{\sum_{i = 1}^{n} \var{y_i} \pdv{f}{y_i'}}_{\alpha}^{\beta}.
\]

Suppose boundary conditions are identical on \(y_i\)s. Then \(\vb{y} \pqty{\vb{x}}\) optimises \(F\) if and only if
\[
    \pdv{f}{y_i} - \dv{x} \pdv{f}{y_i'} = 0
\]
for all \(i = 1, 2, \ldots, n\). This gives \(n\) second-order ODEs.

\subsection{Geodesics of the Euclidean Plane}

We return to \autoref{prob:shortest-path} at the start of the course, and we aim to find that is the shortest curve between two points \(A, B\) in Euclidean space.

The first question we have is how to parameterise this curve. There are two choices.

\paragraph{Parameterise using \(x\).} We write \(y = y\pqty{x}\), for \(x_A \leq x \leq x_B\), subject to constraints \(y\pqty{x_A} = y_A\), and \(y\pqty{x_B} = y_B\). In this case, we have the length functional satisfies
    \[
        L\bqty{y} = \int_{x_A}^{x_B} \sqrt{1 + \pqty{y'}^2} \dd{x}.
    \]

    We have \(f = \sqrt{1 + \pqty{y'}^2}\), and in particular, \(\pdv*{f}{y} = 0\). \autoref{eq:euler-lagrange} implies that \(\pdv{f}{y'}\) is a constant, so
    \[
        \frac{y'}{\sqrt{1 + \pqty{y'}^2}}
    \]
    is constant. Hence, \(y'\) is constant, and \(y = mx + c\) is some straight line.

    Then \(m\) and \(c\) are fixed by the boundary conditions.

But by doing this, we excludes certain types of curves (that does not satisfy the vertical line test). So we consider another parameterisation.

\paragraph{Use arbitrary parameter \(t \in \ccintv{0}{1}\).} We have \(t = 0\) at \(A\), \(t = 1\) at \(B\), and \(x = x\pqty{t}\), \(y = y\pqty{y}\). Then we have
\[
    L\bqty{\vb{x}} = \int_{0}^{1} \sqrt{\dot{\vb{x}}\pqty{t}^2} \dd{t} \qqtext{where} \dot{\vb{x}} = \pqty{\dv{x}{t}, \dv{y}{t}}.
\]

The boundary conditions are \(x\pqty{0} = x_A\), \(x\pqty{1} = x_B\), and similar for \(y\). We have
\[
    f = \sqrt{\dot{x}^2 + \dot{y}^2},
\]
and \autoref{eq:euler-lagrange} on \(x\) imposes that
\[
    \dv{t} \pdv{f}{\dot{x}} - \pdv{f}{x} = 0.
\]

So \(\pdv{f}{\dot{x}}\) is some constant \(c\), and hence
\[
    \frac{\dot{x}}{\sqrt{\dot{x}^2 + \dot{y}^2}} = C
\]
for some constant \(C\). Similarly,
\[
    \frac{\dot{y}}{\sqrt{\dot{x}^2 + \dot{y}^2}} = S
\]
for some constant \(S\). Note \(C^2 + S^2 = 1\), so let \(C = \cos \theta\), \(S = \sin \theta\) for some \(\theta\).

If we switch to use arc length \(l\) as a parameter,
\[
    \dv{l}{t} = \sqrt{\pqty{\dv{x}{t}}^2 + \pqty{\dv{y}{t}}^2} = \sqrt{\dot{x}^2 + \dot{y}^2}.
\]

Then
\[
    \dv{x}{t} = \flatfrac{\dv{x}{t}}{\dv{l}{t}} = \frac{\dot{x}}{\sqrt{\dot{x}}^2 + {\dot{y}}^2} = \cos \theta
\]
and similarly \(\dv*{y}{l} = \sin \theta\).

So \(x = x_A + l \cos \theta\), \(y = y_A + l \sin \theta\) which is a completely arbitrary straight line.

\begin{remark}
    There is a good reason why we start with \(t\) -- since if we use \(l\) initially, then the upper limit of the integration \emph{will depend on the curve}, but our derivation of \autoref{eq:euler-lagrange} assumed that the boundaries are fixed.
\end{remark}